\section{Running It in the Browser}
\label{sec:web}

The activity works on paper and always will. The browser version exists because
it removes three things a facilitator would otherwise carry: keeping every team
on the same phase, adjudicating from memory, and transcribing record sheets
afterwards. If it fails mid-class, keep playing on paper --- every mechanic works
with the physical deck, a printed grid and six coins.

\subsection{Opening It}

Everyone uses the same link. Roles are chosen, not issued.

\begin{center}
\small
\begin{tabularx}{\textwidth}{@{}l l X@{}}
\toprule
\textbf{Who} & \textbf{Opens} & \textbf{Does} \\
\midrule
Facilitator & the link, then \emph{Take the facilitator role} &
Sets scenario, condition, budget and phase; deals; adjudicates; exports \\
Players & the plain link & Join a team, place the threat, list measures, buy
controls, write justifications \\
Projector & add \texttt{\#projector} to the URL &
Read-only shared screen: countdown, every team on one grid \\
\bottomrule
\end{tabularx}
\end{center}

Add \texttt{?session=tue-am} to run more than one class at once. Sessions are
entirely separate.

\textbf{The facilitator code} is chosen when you claim the role. Only someone
with it can take the role while you are active, only you can release it, and it
frees itself about two minutes after your last action so a dead laptop does not
strand the class. Write it down.

\subsection{What the Interface Does For You}

\begin{itemize}
\item \textbf{Phase guidance} appears above each phase: what to do, the one thing
that matters, and what \emph{done} looks like. Teams can collapse it.
\item \textbf{The threat model locks} when you load a scenario, so nobody rerolls
a pre-built one by accident. There is a Lock/Unlock toggle if you want teams
drawing their own.
\item \textbf{Coverage is hidden until Phase 3.} Teams cannot watch a square turn
green while they shop, which is what stops them clicking cards to see what
lights up rather than reasoning.
\item \textbf{Modifier justifications queue} for you with an amber badge showing
how many are waiting.
\item \textbf{The scenario's own debrief questions} appear in your panel at Phase
3, and as separate answer boxes for teams.
\item \textbf{Export} produces one row per team with everything the study scores.
\end{itemize}

\subsection{Accessibility}

The header carries a text-size control (100, 115, 130\%) and a light/dark
toggle. Both save per device, not per session, so one student enlarging text does
not resize the board for the room. Mention this at the start; students rarely
find it on their own and some will not ask.

All colours meet WCAG AA contrast in both themes, and nothing in the interface
depends on colour alone --- the matrix uses shape and letters as well as hue.

\subsection{When Something Looks Wrong}

The build version sits in grey beside \emph{How to play}. If a student reports
odd behaviour, check theirs matches yours; a stale cached page is the most common
cause. A hard reload fixes it.

If a screen goes blank, the app shows a readable error rather than white space.
Have the student read the line to you --- it names the fault.

% =====================================================================
\section{Running the Controlled Study}
\label{sec:study}

Skip this section unless you are collecting research data. Everything in it
exists to keep one comparison honest.

\subsection{The Two Arms}

\begin{center}
\small
\begin{tabularx}{\textwidth}{@{}l X X@{}}
\toprule
& \textbf{Solutions deck} & \textbf{Control} \\
\midrule
Phase 1 & Identical & Identical \\
Phase 2 & List the measures you would recommend, \emph{then} buy controls
against a budget & List the measures you would recommend. No cards, no budget \\
Phase 3 & Identical & Identical \\
Debrief & Identical & Identical \\
Guidance, timing, interface & Identical & Identical \\
\bottomrule
\end{tabularx}
\end{center}

\begin{intervene}{The one thing that would invalidate the study}
The arms must differ \emph{only} in the Solutions cards and the budget. If the
control group gets less guidance, less time, or a different facilitator manner,
you have compared a guided session against an unguided one and any difference you
find means nothing. The software holds guidance, timing and interface constant.
\textbf{You have to hold your own behaviour constant.}

Read the same scenario the same way. Intervene by the same rules. Do not explain
extra because a control team looks stuck.
\end{intervene}

\subsection{Why Both Arms List Measures}

Both conditions ask teams to write down the measures they would recommend, and in
the Solutions arm that list sits \emph{above} the cards and explicitly invites
things they could not afford.

This is not decoration. One of the study's measures counts distinct mitigations
raised per scenario. If that count came from purchases, a six-point budget would
cap the Solutions arm at roughly five while the control arm listed as many as it
liked --- the measure would be biased against the deck by construction. Asking
both arms the same question fixes it, and the budget's effect shows up in the
measures that are about \emph{choosing well} rather than \emph{choosing many}.

Say this to teams in the Solutions arm: \say{List everything you would recommend,
including what you cannot afford. Buying is a separate question.}

\subsection{Setting the Condition}

Set it in the \emph{Condition} panel before anyone joins, or open the session
with \texttt{?condition=control} in the URL. The arm is stored in the session, so
a student cannot switch it by editing their address bar.

\textbf{It locks as soon as any team acts.} Changing arms mid-session would mix
two conditions in one dataset, so the control greys out with the reason shown.

Run the arms as separate sessions:

\begin{center}\small\ttfamily
?session=fall-a\\
?session=fall-b\&condition=control
\end{center}

\subsection{Before, During, After}

\textbf{Before.} Confirm IRB approval covers what you are collecting. Confirm the
build version is current on every device. Decide who reads scenarios aloud and
have them do it for both arms. Turn injects off; they add variance not held
constant across teams.

\textbf{During.} Keep to the timer. Record your interventions by rung on the
facilitator form --- unequal intervention between arms is a confound you can
measure only if you wrote it down. Note anything that broke.

\textbf{After.} Export before clearing anything; neither reset archives for you.
Complete the facilitator form the same day, while you still remember which team
argued and which went quiet.

\begin{note}{What the export gives you}
One row per team: condition, scenario, threat model, both matrix placements,
budget, spend, purchases, everything ever purchased, accepted modifiers, the
measures list and its count, residual risk, each debrief answer separately, and
the closing recommendation. That is most of the study's data without a single
record sheet being transcribed.

It does not capture what teams \emph{said}. If you want the argument as well as
the outcome, you still need an observer or a recording, and consent for it.
\end{note}
